\section{Discussion}
\label{sec:discussion}

The evaluation indicates that TPS-ZKPI reduces the policy-dependent signals in normalized public-chain artifacts below our multi-seed leakage criterion.
Across all PCH, PI, and SPP tasks, supervised classifiers stay below that criterion, and unsupervised clustering shows no recoverable structure: ARI stays within $\pm 0.002$ of chance on every target, and the one elevated NMI (ICD-10) is consistent with the finite-sample inflation quantified by the null simulation of Section~\ref{sec:eval-nmi-sensitivity}.
The gas-padding ablation (Section~\ref{sec:eval-ablation}) sharpens this conclusion: the same attacks become measurably effective once gas padding is removed.
This attributes the gas-channel protection specifically to randomized gas padding.
Content hiding itself rests on the commitment-only interface and the content-independent transcript normalization, with padding additionally suppressing the residual gas-channel signal that is otherwise recoverable for some content labels.
Together with fixed-length policy representation, reusable interpretation, and commitment
binding, these mechanisms limit the extent to which public-chain observers infer hidden policy content or structural properties from authorization transactions.

These results support a scoped privacy claim.
The evaluation focuses on a metadata-only public-chain adversary that observes commitments, proofs, public parameters, calldata, events, and gas-related artifacts, but not TPS-side witness construction or off-chain policy material.
Within this scope, the residual cost is the randomized padding's gas overhead: about $38\%$ on access gas and $40\%$ on commit gas (Section~\ref{sec:eval-ablation}).

Boundary cases above the $0.55$ AUC line do occur---PCH chapters ch6 and ch7, and the PI \texttt{semantic\_equivalent} pairs---but none reproduces across seeds, and we report them. Stronger adversaries---compromised proof services, adaptive requester-side probing, or off-chain side channels---remain outside this scope and motivate the reduced-trust extensions discussed below.

The TPS trust assumption is the main architectural tradeoff.
Moving confidentiality-sensitive proof generation to the TPS keeps policy-bound proving information away from requesters and public smart contracts while preserving on-chain verifiability, but it shifts part of the confidentiality burden to the service layer; the provider-operated TPS used here concentrates that burden at the provider rather than at an independent service (Section~\ref{sec:problem}).
This tradeoff leaves on-chain verifiability intact but does affect liveness and accountability: the TPS or provider can refuse to produce a proof, and such a refusal may be indistinguishable from a legitimate denial at the on-chain interface.

Attribute origin is only partially delegated to the TPS.
For the profile-backed DUO terms, the circuit verifies the EdDSA-Poseidon signature described in Section~\ref{sec:design} and derives eight policy attributes from the seven signed claims.
A compromised TPS cannot modify the signed claims, policy commitment, requester pseudonym, resource, validity interval, or credential identifier without invalidating the signature.
The implemented relation rejects a non-committed issuer, an invalid signature under the committed key, an expired or premature credential, and substitutions of any signed field.
Other evidence has a narrower meaning: ontology proofs establish that a supplied code belongs to a provider-committed set, and requester allowlists establish that a holder pseudonym belongs to a provider-committed list; they do not authenticate every underlying real-world fact.
Moreover, because the TPS receives the holder secret and credential witness in the current workflow, holder binding does not protect against a malicious TPS reusing otherwise valid evidence within its signed scope.

The remaining trust anchor is therefore issuer governance.
The prototype binds the proof to a policy-selected issuer-key commitment, but it has no independent registry establishing that this key belongs to an accredited healthcare authority, and it does not check credential revocation.
The evaluation uses a deterministic test-only issuer key to exercise the complete signature path; it is not evidence of institutional issuance or production key management.

A concrete deployment can close this gap without changing the Boolean interpreter.
First, a governed issuer registry would publish schema-qualified issuer keys and their validity epochs; registration would validate a canonical, subgroup-safe public key and require proof of possession.
Second, an issuer would sign a domain-separated credential containing a holder-key commitment, claim vector, authorization scope, validity interval, and unique credential identifier.
Third, the circuit would verify both the credential signature and a Merkle proof that the issuer key is active in the registry, while continuing to derive policy attributes from the signed claims.
The holder would retain its secret key and authorize each request by signing the context digest and request epoch; verifying this second signature inside the circuit would prevent a TPS that observes one credential presentation from generating proofs for fresh sessions.
Finally, an epoch-specific sparse-Merkle or accumulator root would support a private non-revocation witness for the credential identifier.
The registry and revocation roots can occupy fixed-format public or committed parameter slots, preserving the policy-independent verification interface.
If an existing healthcare credential uses a non-circuit-friendly signature such as Ed25519 or ECDSA, integration requires either verifying that signature in-circuit at higher constraint cost or having the issuer additionally issue a SNARK-friendly credential; re-signing by the TPS would not preserve issuer-origin assurance.

As described in Section~\ref{sec:design}, the policy commitment covers both the padded encoding $P_{\mathrm{enc}}$ and its active length $\ell$.
Because $\ell$ also determines how many tokens the interpreter evaluates, an authorizing proof for $h$ cannot be based on a shorter valid prefix.
This integrity guarantee is separate from our privacy evaluation, which measures whether public-chain artifacts reveal information about the hidden policy.

The authorization context is likewise part of the proved relation.
The circuit constrains the public digest $d_{\chi}$ to the requester identifier, resource, and nonce used in the witness, and the ACM recomputes the same digest from the transaction parameters.
Consequently, changing any of these values invalidates the request.


The guarantee is scoped to the bound tuple within one ACM deployment.
The digest does not include the transaction submitter or a chain/contract domain, so exact-context front-running and cross-deployment reuse require separate deployment-level controls.
Moreover, $\mathit{id}_{\mathrm{req}}$ is a pseudonymous authorization identifier and is not, by itself, proof of a real-world identity or of the origin of every private attribute.

Three directions follow.
First, adding an explicit chain/contract domain and an authenticated submission channel would extend context binding beyond a single deployment and reduce front-running risk.
Second, accountable proof obligations---for example, on-chain commitments to serve requests together with timeout-with-evidence mechanisms---would make proof refusals attributable.
Third, the reduced-trust TPS instantiations of Section~\ref{sec:problem} (collaborative-SNARK or TEE-backed proving) would relax the trust concentration itself, including the attribute-origin limitation; realizing and evaluating them is left to future work.
